% En esta seccion se expresan sus conclusiones del trabajo realizado de manera concreta y coherente.
\section{Conclusiones}

    \subsection{Basilio Illescas Andrés}
    Los objetivos de esta entrega se cumplieron: las tres primeras actividades del manual quedaron implementadas y verificadas, con el plano de $2\times2$, el cubo de $2\times2\times2$ y el círculo de radio unitario centrados en el origen y observables desde las dos perspectivas asignadas a \texttt{F1} y \texttt{F2}. Lo que realmente me llevo es que el modelado no consiste en escribir coordenadas sino en elegir una descomposición: el plano sólo pedía leer bien el constructor, el cubo pedía verificar por caras que no faltara ninguna, y el círculo pedía decidir cómo se aproxima una curva y con cuántos lados, de modo que la figura dejó de escribirse y pasó a generarse. La lección de método vino de los dos errores de la clase \texttt{Circle} del proyecto base: el de los colores se veía de inmediato, pero el índice fuera de rango no produjo ningún síntoma en pantalla y sólo apareció al imprimir los arreglos, lo que me dejó claro que en gráficos «se ve bien» no equivale a «está bien» y que por eso valió la pena agregar la tecla que alterna entre relleno y alambre.

    \subsection{Cruz Macedo Samuel Santiago}
    Al principio esta práctica me pareció más sencilla que la anterior, porque ya no había que pelearse con tres matrices, pero la dificultad simplemente se movió de lugar: en la Práctica 02 el problema era decidir desde dónde se mira la figura y aquí es decidir de qué está hecha. Lo que más me sirvió fue entender la diferencia entre la primitiva que uno imagina y la que la tarjeta dibuja, porque no existe una instrucción para dibujar un círculo: son setenta y dos triángulos que comparten el centro, y al bajar el número de lados a doce la figura se volvió un polígono evidente, lo que deja ver que toda curva en gráficos es una aproximación. De la clase \texttt{Circle} del proyecto base me quedó una lección aparte, ya que los tres sectores de color plano que salían no eran culpa del shader sino de que \texttt{Vertex} no se inicializa sola: fue la primera vez que vi un programa correr sin marcar error y aun así entregar algo indefinido. Considero que los objetivos se cumplieron, y me queda claro que esto es la base de lo que sigue, porque el cilindro y el cono son este mismo círculo repetido o unido a un punto.

    \subsection{Medel Piña Ángel de Jesús}
    Con esta práctica pude comprender cómo se construye una escena a partir de figuras básicas y cómo las transformaciones geométricas permiten reutilizar una misma geometría varias veces. Al armar la escena noté que el orden en que se aplican la traslación, la rotación y el escalamiento sí importa, ya que al multiplicarse como $T\,R\,S$ primero se escala la figura en sus propios ejes, luego se orienta y al final se coloca en su posición. Gracias a esto pude, por ejemplo, recortar el cilindro y el cono en el eje $z$ para ajustar su altura antes de rotarlos y dejarlos de pie sobre el piso, además de calcular la altura de cada objeto para que quedara apoyado sobre la superficie en lugar de atravesarla o quedar flotando.

    También me di cuenta de que el color de las figuras no es una propiedad que se pueda cambiar al momento de dibujarlas, sino que forma parte de los datos de cada vértice que se envían al \textit{buffer} al crear la malla. Por eso, para tener objetos del mismo tipo con colores distintos fue necesario construir una malla por cada color, mientras que para repetir un objeto con el mismo color bastaba con dibujar la misma malla con diferentes matrices de modelo. Por último, implementar las caras de la esfera me ayudó a entender cómo una rejilla de vértices guardada en un solo arreglo se puede convertir en triángulos calculando directamente sus índices, y cómo el operador módulo permite cerrar la figura sin dejar costuras abiertas.

    \subsection{Perez Olvera Alexis Abraham}
    Esta práctica me dejó claro que la parte difícil de construir una geometría 3D no son las fórmulas matemáticas, que casi siempre ya venían resueltas en el proyecto base, sino los detalles que uno da por sentados: inicializar bien una estructura, cerrar un contorno, no repetir un vértice de más. El caso del índice fuera de rango en el círculo, que no producía ningún error visible en pantalla, fue el que más me hizo entender que una figura puede compilar y verse bien sin estar realmente correcta, y que conviene desconfiar un poco de lo que se ve bien a simple vista.

    También me quedó una intuición más clara de algo que antes usaba sin pensarlo demasiado: cualquier modelo 3D, por sofisticado que parezca, es en el fondo lo mismo que hice aquí a mano —vértices, colores y triángulos que los conectan. Verlo funcionar, con sus errores incluidos, terminó explicándome esa idea mejor que cualquier lectura previa.

    
\subsection{Segura Cedeño Luisa María}

Esta práctica me ayudó a entender cómo se construyen figuras tridimensionales más complejas a partir de triángulos básicos. Al principio pensé que bastaba con usar el archivo del cono que ya venía en el proyecto base, pero al correrlo me di cuenta de que tenía varios problemas: casi todo se veía negro porque le faltaba ponerle color a la mayoría de los puntos, la punta de arriba estaba repetida muchas veces y la figura quedaba hueca por no tener tapa abajo.

Al ver el trabajo completo con el plano, el cubo, el círculo, la esfera y el cono, entendí que toda geometría tridimensional, sin importar qué tan redonda parezca, termina siendo una descomposición en triángulos y vértices ordenados. Además, poder alternar entre las vistasy usar el modo de alambre fue indispensable para comprobar que ninguna cara quedara mal conectada, demostrando que en gráficos por computadora no basta con que algo se vea bien a simple vista, sino que su estructura interna debe estar bien construida.

